% Chapter 5: Thinking Together
\chapter{Thinking Together}
\label{ch:thinking-partner}

\begin{itshape}
Your prompts are precise, your CLAUDE.md is configured, and your tests pass.
But every interaction follows the same pattern: you delegate, Claude executes,
you verify. Something is missing---you are using Claude as a tool you configure,
not a collaborator you think with.
\end{itshape}

The previous chapters answer ``how do I tell Claude what to do?''
This chapter answers a different question: how do I \emph{think together}
with Claude?

The core metaphor: Claude is a thinking mirror. It reflects your codebase
back to you. Where the reflection is clear, your code is clear.
Where it is distorted---where Claude misunderstands, guesses, or asks
for clarification---you have found a documentation gap, a hidden assumption,
or an unclear interface.

\begin{warnbox}[An Honest Caveat]
Claude is a thinking partner with no memory, a tendency to agree,
and occasional confident wrongness. The techniques in this chapter
work \emph{because} they structure the collaboration to counteract
these weaknesses---not despite them.
\end{warnbox}

\begin{whybox}[Why This Chapter Exists]
Every chapter so far answers ``how do I configure Claude to do what I want?''
This one answers ``how do I use Claude to think more clearly?''
The shift is from delegation to collaboration---from configure-delegate-verify
to think-together-discover-build-better.
\end{whybox}

% -------------------------------------------------------------------
\section{Hypothesis-Driven Debugging}
\label{sec:hypothesis-debugging}

When a bug appears, the default instinct is to paste the traceback and say
``fix this.'' Claude will try---and often succeed on the first attempt.
But for deeper bugs, this approach produces patches that address symptoms
rather than root causes.

\practitioner{Structure debugging as hypothesis testing, not error correction.
Claude is a pair programmer who catches different errors because it has
no sunk cost in your approach.}

\begin{fullwidth}
\begin{beforeafter}[Debugging]
\before
\begin{minted}{text}
Fix this error:
  File "feature_pipeline.py", line 47
  ValueError: negative value in non-negative
  feature 'days_since_last_purchase'
\end{minted}

\after
\begin{minted}{text}
I see a ValueError in feature_pipeline.py
line 47: negative values in a feature that
should be non-negative. Three hypotheses:

1. Log transform applied before clipping
   negative deltas
2. Currency conversion introduces negatives
   for returns/refunds
3. Timezone mismatch causes date subtraction
   overflow

For hypothesis 1: write a test with a purchase
  date after the reference date.
For hypothesis 2: write a test with a refund
  transaction in the events table.
For hypothesis 3: check if timestamps are
  timezone-aware or naive.

Run hypothesis 1 first---it is the most
likely given the stack trace.
\end{minted}
\end{beforeafter}
\end{fullwidth}

\margintip{``Three hypotheses for this failure. Design a minimal test
for each.''}

The second prompt is longer but produces better outcomes: a root-cause
diagnosis instead of a symptom patch. Each hypothesis becomes a testable
prediction. When the test confirms one, you fix the cause, not the effect.

\margintip[Tip]{``Run this one test with this one input. Does it pass?''
Isolation is the fastest path to root cause.}

\keyinsight{Rubber duck debugging works because articulating the problem
reveals the gap. Claude-as-duck works better because the duck talks back---it
asks questions you would not ask yourself, because you already ``know''
the answer. (Sometimes you are wrong about what you know.)}

\begin{trybox}
Pick a recent data pipeline bug---ideally one where results looked plausible
but were wrong (e.g., feature drift, silent NaN propagation, leakage).
Write three hypotheses for what caused it.
Ask Claude to design a minimal test for each hypothesis.
Would this approach have found the root cause faster?
\end{trybox}

% -------------------------------------------------------------------
\section{Tests as Thinking Tools}
\label{sec:tests-as-thinking}

In \cref{ch:testing}, tests serve verification: does the code do what it claims?
Here, tests serve a different purpose---exploration:
\emph{what should the code do at all?}

\begin{beforeafter}[Tests: Verification vs.\ Exploration]
\before
``Write tests for this function.''

\after
``I am not sure what should happen at the boundary.
Write 5 test cases exploring: empty input, single element,
duplicates, negatives, overflow.
Which behaviors surprise me?''
\end{beforeafter}

\begin{whybox}[Why Tests Before Requirements]
Writing test cases forces you to articulate expectations you have not
stated. When a test case surprises you---the function does something
you did not expect---you have discovered a requirement that was implicit.
The test did not verify your code. It interrogated your assumptions.
\end{whybox}

\margintip{``What invariant should always hold for this function,
regardless of input?'' This question surfaces design decisions
hiding as implementation details.}

\subsection{Property-Based Tests as Assumption Detectors}

\practitioner{Property-based tests are particularly powerful as thinking tools.
Instead of testing specific inputs, they test \emph{invariants}---properties
that must hold for all inputs.}

\begin{minted}{text}
Prompt: "What properties should always be true
for this sort function?
- Output length == input length
- Output is ordered
- Output is a permutation of input
Write a property-based test for each."
\end{minted}

When a property-based test fails on a randomly generated input, it has found
an edge case neither you nor Claude anticipated. That failure is more
valuable than a passing test---it reveals a boundary in your understanding.

\marginnote[Cross-Ref]{For test infrastructure and verification criteria,
see \cref{ch:testing}. For characterization tests, see \cref{ch:projects}.}

% -------------------------------------------------------------------
\section{Surfacing Hidden Assumptions}
\label{sec:hidden-assumptions}

Every system rests on assumptions---about scale, about usage patterns,
about what will never change. Most are invisible until they break.
Claude can surface them, not because it is smarter, but because it has
no investment in your preferred approach.

\begin{beforeafter}[Architecture Review]
\before
``Review my database schema.''

\after
``Here is my feature store schema. I designed it assuming:
(1) features are computed daily in batch, not real-time,
(2) training and serving use the same feature computation,
(3) feature drift is monitored externally.
Which assumption is most likely wrong in 6 months,
and what breaks when it does?''
\end{beforeafter}

The second prompt names the assumptions explicitly, then asks Claude
to stress-test them. This is not code review---it is assumption review.

\begin{tipbox}[The Pre-Mortem Prompt]
``Imagine this feature has failed in production six months from now.
What are the three most likely causes?
Work backwards from failure to the design flaw that enabled it.''

Pre-mortems are more effective than post-mortems because they cost nothing
and can change the design. Ask Claude to run one before every
architecture decision.
\end{tipbox}

\begin{tipbox}[The Feynman Prompt]
``Explain my auth flow as if I just joined the team and need to modify it.
Where did you have to guess because the code does not make intent clear?''

Gaps in Claude's explanation are gaps in your documentation.
What Claude cannot explain, a new hire cannot understand.
\end{tipbox}

\margintip{``What am I assuming that I have not written down?''
The most dangerous assumptions are the ones nobody states.}

% -------------------------------------------------------------------
\section{Quick Wins That Compound}
\label{sec:quick-wins}

These five techniques take minutes to apply and improve every subsequent
Claude interaction. The metaphor: Claude reads your codebase like a
developer who just joined the team. Good onboarding materials help
Claude the same way they help humans.

\subsection{Docstrings as Triple-Duty Documentation}

A docstring serves three audiences simultaneously:
you (in six months), Claude (right now), and your teammates (always).

\begin{fullwidth}
\begin{beforeafter}[Docstrings]
\before
\begin{minted}{python}
def calculate_risk(portfolio, window):
    scores = []
    for asset in portfolio:
        s = _score(asset, window)
        if s is not None:
            scores.append(s)
    return np.mean(scores) if scores else 0.0
\end{minted}

\after
\begin{minted}{python}
def calculate_risk(
    portfolio: list[Asset],
    window: int = 30,
) -> float:
    """Average risk score across portfolio.

    Args:
        portfolio: Assets to evaluate.
            Empty portfolio returns 0.0.
        window: Lookback days. Default 30.

    Returns:
        Mean risk score in [0, 1].
        Returns 0.0 if no assets scorable.

    Note:
        Skips assets where _score() returns
        None (insufficient history).
        This is intentional, not an error.
    """
\end{minted}
\end{beforeafter}
\end{fullwidth}

The ``Note'' section is the highest-value addition: it explains a design
decision that would otherwise look like a bug. Claude will not
misinterpret the \code{None} handling because the intent is documented.

\subsection{Type Hints as Contracts}

Five seconds to write, five minutes of debugging saved.
\code{window: int = 30} tells Claude the type, the default,
and the name---three constraints in five characters.
Without the hint, Claude may pass a string, a float, or a \code{timedelta}.

\subsection{Error Messages That Debug Themselves}

This connects directly to Principle~4 from \cref{sec:four-principles}:
fail fast with diagnostics. An error message that includes what went wrong,
what was expected, and what to do about it saves Claude the same
investigation time it saves you.

\subsection{Code Archaeology}

For brownfield work, Claude makes an excellent excavation partner:

\begin{minted}{text}
Prompt: "Why might the original author have
written this as a nested loop instead of a join?
What constraint explains this design choice?"
\end{minted}

Instead of assuming the code is wrong, this prompt assumes the code is
\emph{explained by something you do not yet see}. Claude often finds the
constraint---a database limitation, a legacy API, a performance requirement---
that makes the original design rational.

\subsection{README-Driven Development}

Write the README first. Then ask Claude:

\begin{minted}{text}
Prompt: "Read this README. What questions does a
new developer still have after reading it?"
\end{minted}

\margintip{``Read this function. What would you need to know
before modifying it?'' Gaps in the answer are gaps in your code.}

\begin{trybox}
Pick one function Claude frequently reads in your project.
Add a docstring with Args, Returns, and one Note explaining a non-obvious
design decision. Compare Claude's output on a task involving that function
before and after the docstring.
\end{trybox}

% -------------------------------------------------------------------
\section{Architecture Decision Records}
\label{sec:adr}

When you make an architecture decision \emph{with} Claude,
the conversation captures not just what you decided but why,
and what alternatives were considered. An Architecture Decision Record
(ADR) preserves this reasoning for your future self.

\begin{minted}{markdown}
# ADR-007: Offline Feature Computation

## Context
Feature computation runs in nightly batch.
Some features are stale by 12 hours at serving time.

## Decision
Keep batch for training features.
Add streaming for 3 real-time features
(last login recency, cart value, session count).

## Alternatives Considered
1. All streaming (rejected: 10x infrastructure cost)
2. Faster batch, hourly (rejected: still stale)
3. Feature caching with TTL (rejected: cache
   invalidation complexity)

## Consequences
- Two feature computation paths to maintain
- Real-time features need monitoring for drift
- Training/serving skew possible for 3 features

## Assumptions to Revisit
- 3 real-time features sufficient for next 6 months
- Streaming infra handles peak load (Black Friday)
- Feature drift alerts catch training/serving skew
\end{minted}

\keyinsight{The most valuable ADR section is ``Alternatives Considered.''
Claude suggests alternatives you would not---not because it is smarter,
but because it has no investment in your preferred approach.
A human colleague might hesitate to challenge your solution.
Claude does not hesitate.}

\margintip{``Draft an ADR for this decision. What alternatives
should I have considered?'' Enable extended thinking
(\cref{sec:thinking}) when drafting ADRs.}

\marginnote[Cross-Ref]{For plan documents and large-task workflows,
see \cref{sec:plans}.}

% -------------------------------------------------------------------
\section{Quick Reference}
\label{sec:thinking-reference}

\begin{itemize}
  \item Claude is a thinking mirror: distortions reveal gaps in your code
  \item Debugging: three hypotheses, minimal tests, isolate blame
  \item Tests as exploration: discover requirements, not just verify them
  \item Assumption surfacing: name your assumptions, then stress-test them
  \item Pre-mortem: imagine failure, work backwards to design flaw
  \item Docstrings, type hints, error messages: minutes to add, compound forever
  \item Code archaeology: assume intent before assuming error
  \item ADRs: capture alternatives considered, not just the decision
\end{itemize}

\subsection{Building a Prompt Pattern Library}

The techniques in this chapter can be codified as team-shared resources.
The hypothesis debugging workflow becomes a custom command.
The pre-mortem prompt becomes a line in your shared CLAUDE.md.
The ADR template becomes a skill.

\practitioner{Start personal, then share. Once a prompt pattern saves you
time three sessions in a row, promote it from your terminal history
to a command or skill that the whole team can use.}

\marginnote[Cross-Ref]{For team-shared patterns and governance,
see \cref{ch:team}.}

% -------------------------------------------------------------------
\begin{trybox}
Choose one technique from this chapter and apply it today:
\begin{enumerate}
  \item \textbf{Easiest}: Add a docstring to one feature engineering function
    Claude reads often. Include the expected DataFrame schema.
  \item \textbf{Most impactful}: Use the hypothesis debugging workflow
    on your next pipeline bug.
  \item \textbf{Most strategic}: Draft an ADR for an upcoming
    feature store or model serving decision.
\end{enumerate}
Pick one. Do it now. The techniques compound---but only if you start.
\end{trybox}
